
\subsection{Definitions and Notations}



A \emph{signature} of arity $n$ is a function
$f:\{0,1\}^n\to\mathbb C$. 
%We also write $\operatorname{arity}(f)=n$ and $f^\alpha=f(\alpha)$.  
%Unless a row orientation is explicitly indicated, the truth table of $f$ is the column vector in lexicographic order.  
For a bit string $\alpha=\alpha_1\cdots\alpha_n$, let $|\alpha|=n$,
\(
\operatorname{wt}(\alpha)=\sum_{j=1}^n\alpha_j,
\)
and denote its bitwise complement by $\overline{\alpha}$.  
The symbols
$0^n=\vec 0^{\,n}$ and $1^n=\vec 1^{\,n}$ denote the all-zero and all-one
strings; the superscript is omitted when its value is clear.  Addition
$\oplus$ of bit strings is coordinatewise over $\mathbb Z_2$.
The \emph{support} of $f$ is
\(
\mathscr S(f)=\{\alpha\in\{0,1\}^n:f(\alpha)\ne0\}.
\)
The signature is zero, written $f\equiv0$, if $\mathscr S(f)=\emptyset$.
A signature is \emph{symmetric} if its value depends only on Hamming weight.
For such a signature of arity $n$, we use the compressed notation
\(
f=[f_0,f_1,\ldots,f_n],
\)
where $f_j$ is the value on every input of weight $j$. 
Let
\(
\Delta_0=[1,0], \,\Delta_1=[0,1].
\)
For $n\ge1$, the equality signature $(=_n)$ is one on $0^n$ and $1^n$ and
zero elsewhere.  The binary disequality signature is
$(\neq_2)=[0,1,0]$. 
More generally, $\neq_{2n}$ denotes the indicator of
\(
x_1= x_2=\ldots= x_n\neq
x_{n+1}=x_{n+2}=\ldots= x_{2n}.
\)
We
write $\mathcal{EQ}=\{=_n\mid n\ge1\}$ and
$\mathcal{DEQ}=\{\neq_{2n}\mid n\ge1\}$.
Also write $\mathcal{EQ}_k=\{=_{nk}\mid n\ge1\}$.
Let
\(
\mathscr E_n=\{\alpha\in\{0,1\}^n:\operatorname{wt}(\alpha)
\text{ is even}\},\,
\mathscr O_n=\{\alpha\in\{0,1\}^n:\operatorname{wt}(\alpha)
\text{ is odd}\}.
\)
A signature has \emph{even parity} or \emph{odd parity} if its support is
contained in $\mathscr E_n$ or $\mathscr O_n$, respectively; it has
\emph{parity} in either case.

We use $=_2^-$ to denote the binary signature $(1, 0, 0, -1)$ and $\neq_2^-$ to denote the binary signature $(0, 1, -1, 0)$.
We may also write $=_2$ as $=_2^+$ and $\neq_2$ as $\neq_2^+$. 
%We use $=_2^{\pm}$ to denote $=_2^+$ or $=_2^-$, and use $\neq_2^{\pm}$ to denote $\neq_2^+$ or $\neq_2^-$.
Let $\mathcal{B}=\{=^+_2, =_2^-, \neq_2^+, \neq_2^-\}$.
We call them Bell signatures which correspond to Bell states $|\Phi^+\rangle=|00\rangle+|11\rangle$, $|\Phi^-\rangle=|00\rangle-|11\rangle$, 
$|\Psi^+\rangle=|01\rangle+|10\rangle$ and 
$|\Psi^-\rangle=|01\rangle-|10\rangle$ in quantum information science \cite{bell1964einstein}.

We use $f_{i_1i_2\ldots i_k}^{a_1a_2\ldots a_k}$ to denote the signature obtained from $f$ by pinning the variables $x_{i_1},x_{i_2},\ldots, x_{i_k}$ to $a_{1},a_2,\ldots, a_k$ respectively.
For example, $f_{i}^a=f(x_1,\ldots,x_{i-1},a,x_{i+1},\ldots,x_n).$
For $S\subseteq[n]$, let $M_S(f)$ be the
$2^{|S|}\times2^{n-|S|}$ matrix whose rows are indexed by assignments
to the variables in $S$ and whose columns are indexed by assignments to
the remaining variables, both in lexicographic order.  We also write
$M_{S,T}(f)$ when the ordered row and column variable lists $S,T$
partition the variables.  In particular,
\[
M_i(f)=
\begin{bmatrix}\mathbf f_i^0\\ \mathbf f_i^1\end{bmatrix},
\qquad
M_{ij}(f)=
\begin{bmatrix}
\mathbf f_{ij}^{00}\\ \mathbf f_{ij}^{01}\\
\mathbf f_{ij}^{10}\\ \mathbf f_{ij}^{11}
\end{bmatrix},
\]\
where $ {\bf {f}}^{a}_{i}$ denotes the row vector indexed by $x_i=a$ in $M_{i}(f)$, and ${\bf {f}}^{ab}_{ij}$ denotes the row vector indexed by $(x_i, x_j)=(a, b)$ in $M_{i j}(f)$.
For a binary signature $b=b(x_1,x_2)$, $M(b)=\left[\begin{smallmatrix}
    b(0,0) & b(0,1)\\
    b(1,0) & b(1,1)
\end{smallmatrix}\right]$ denotes its $2\times2$ signature matrix.




A \emph{signature grid} over $\mathcal F$ is a pair
$\Omega=(G,\pi)$, where $G=(V,E)$ is a graph and $\pi$ assigns to every
$v\in V$ a signature $f_v\in\mathcal F$ of arity $\deg(v)$, and labels the incident edges $E(v)$ at $v$ with input variables of $f_v$.
We consider all 0-1 edge assignments $\sigma$, and each gives an evaluation $\prod_{v\in V}f_v(\sigma|_{E(v)})$, where $\sigma|_{E(v)}$ denotes the restriction of $\sigma$ to $E(v)$.

\begin{definition}[Holant Problems]
The input to the problem $\Holant(\mathcal{F})$ is a signature grid $\Omega=(G,\pi)$ over $\mathcal{F}$.
The output is the partition function
\[
\operatorname{Holant}({\Omega})
=\sum_{\sigma:E(G)\to\{0,1\}}
 \prod_{v\in V}f_v(\sigma|_{E(v)}).
\]
Bipartite Holant problems $\holant{\mathcal{F}}{\mathcal{G}}$
are {Holant} problems
 over bipartite graphs $H = (U,V,E)$,
where each vertex in $U$ or $V$ is labeled by a signature in $\mathcal{F}$ or $\mathcal{G}$ respectively.
When $\{f\}$ is a singleton set, we write $\Holant(\{f\})$ as $\Holant(f)$ and  $\Holant(\{f\}\cup \mathcal{F})$ as $\Holant(f, \mathcal{F}).$ 
\end{definition}

We say a signature $f$ is \emph{realizable} from $\F$, if $\Holant(\F,f)\leq_T \Holant(\F)$.
The \emph{counting constraint satisfaction} problem \#CSP($\mathcal{F}$) is defined as $\Holant(\mathcal{EQ}\mid \mathcal{F}).$
The problem \#$\mathrm{CSP}_k(\mathcal{F})$ is defined as $\holant{\mathcal{EQ}_k}{\mathcal{F}}$.

For a signature $f$ of arity $n$, define its conjugation $\overline{f}$ by $\overline{f}(\alpha)=\overline{f(\alpha)},\,\forall\alpha\in \{0,1\}^n$.
We define the \emph{arrow reversal} of $f$ by
\(
f^{\textsc{ar}}(\alpha)
:=\overline{f(\overline\alpha)},\,\forall\alpha\in \{0,1\}^n.
\)
We say $f$ satisfies \emph{arrow reversal symmetry} (\textsc{ars}) if
$f=f^{\textsc{ar}}$.  
A signature set $\mathcal F$ is
\emph{conjugate closed}, written $\mathcal F=\overline{\mathcal F}$, if
$f\in\mathcal F$ implies $\overline f\in\mathcal F$.  
In this paper, we focus on the case where $\mathcal{F}$ is conjugate closed.
A set $\mathcal G$
is \emph{arrow-reversal closed}, written
$\mathcal G=\mathcal G^{\textsc{ar}}$, if $g\in\mathcal G$ implies
$g^{\textsc{ar}}\in\mathcal G$.  


Let $\mathcal U$ be the set consisting of the binary zero signature and
all binary signatures $b$ for which
\(
M(b)M(b)^\dagger=\lambda I_2
\text{ for some }\lambda>0.
\)
Similarly, $\mathcal O$ consists of the binary zero signature and the
real-valued binary signatures $b$ satisfying
$M(b)M(b)^{\mathtt T}=\lambda I_2$ for some $\lambda>0$.  
%We set $\widehat{\mathcal O}=K^{-1}\mathcal O$ and $\widehat{\mathcal U}=K^{-1}\mathcal U$. 

We use projective binary matrices in the arity-six argument.  For nonzero
matrices $A,B$, write $A\sim B$ if $A=\lambda B$ for some
$\lambda\in\mathbb C^\times$, and let $[A]$ denote the projective class.
Set
\[
X=N_2,\qquad
Y=\begin{bmatrix}0&-\mathfrak i\\ \mathfrak i&0\end{bmatrix},\qquad
Z=\begin{bmatrix}1&0\\0&-1\end{bmatrix},\qquad
J=\mathfrak iY=\begin{bmatrix}0&1\\-1&0\end{bmatrix}.
\]
Notice that $I_2,X,Z,J$ are matrices of the signatures $=_2^{+},\neq_2^{+},=_2^{-},\neq_2^{-}$, respectively.
For $\varepsilon_1,\varepsilon_2,\varepsilon_3\in\{\pm1\}$, put
\[
C_{\varepsilon_1,\varepsilon_2,\varepsilon_3}
=\frac12\bigl(I_2+\varepsilon_1\mathfrak iX
                 +\varepsilon_2J+\varepsilon_3\mathfrak iZ\bigr).
\]
The projective tetrahedral group and its normal Klein-four subgroup are
\[
\Gamma
=\{[I_2],[X],[Z],[J]\}
 \cup
 \{[C_{\varepsilon_1,\varepsilon_2,\varepsilon_3}]:
        \varepsilon_1,\varepsilon_2,\varepsilon_3\in\{\pm1\}\},
\qquad
K_4=\{[I_2],[X],[Z],[J]\}.
\]
Thus $\Gamma\cong A_4$ and $K_4\cong V_4$, where $A_4$ is the alternating group of order 12 and $V_4$ is the Klein group.
We suppress the projective
brackets when a displayed unitary representative is intended. 
In this paragraph, and everywhere in the quantum arguments, $Z$ is the Pauli
matrix rather than a holographic transformation.
